\begin{textsample}{POS Dim 3 – ro – Score 19.00 – ro/ro\_inf\_31.txt}  \label{ex:f3_pos_127}
Contextos de aprendizagem e desenvolvimento Para propor intencionalidades educativas significativas, docentes da primeira infância podem : Pensando nos \textbf{Bebês} - Oportunizar vivências com \textbf{sons} e músicas por meio de movimento corporal, ou \textbf{batendo}, sacudindo, chacoalhando etc., objetos sonoros diversos. - Organizar exploração das qualidades sonoras de objetos e instrumentos musicais diversos, como sinos, flautas, apitos, coquinhos. - Oportunizar \textbf{brincadeiras} com as possibilidades expressivas da própria voz. - Possibilitar a utilização a seu modo de materiais como \textbf{tintas} caseiras, guache, aquarela etc. na produção visual, ampliando suas possibilidades de exploração da cor. - Possibilitar exploração de materiais gráficos na criação de garatujas e outras formas de expressão.
Pensando nas \textbf{crianças} bem \textbf{pequenas} e \textbf{pequenas} - Oportunizar que cantem sozinhas ou em grupo partes ou frases das canções que já conhecem. - Possibilitar a participação em \textbf{brincadeiras} de roda e jogos musicais. - Organizar e mediar a construção de diferentes objetos sonoros e instrumentos musicais. - Oportunizar que identifiquem os \textbf{sons} da natureza ( cantos de pássaros, “ vozes ” de animais, barulho do vento, da chuva etc. ), da cultura ( vozes humanas, \textbf{sons} de instrumentos musicais, de máquinas, produzidos por objetos e outras fontes sonoras ) e o silêncio. - Oportunizar diálogos sobre as qualidades dos \textbf{sons} de determinados objetos sonoros e instrumentos musicais, ainda que não saibam nomeá -los convencionalmente. - Possibilitar que \textbf{demonstrem} sua preferência por determinadas músicas instrumentais e diferentes expressões da cultura musical brasileira e de outras culturas : canções, acalantos, cantigas de roda, parlendas, trava-línguas etc. - Oportunizar a exploração de diferentes maneiras de produzir \textbf{sons} com o próprio corpo ( \textbf{sons} com as mãos, \textbf{pés} e outras partes do corpo ). - Possibilitar a exploração das relações de \textbf{peso}, tamanho, volume e direção na criação de formas tridimensionais usando diferentes materiais e ferramentas. - Organizar e oportunizar a expressão de sensações conforme exploram objetos ou materiais com \textbf{texturas} diversas. - Possibilitar a criação de formas planas e com volume por meio da escultura, modelagem com barro e argila, etc. - Oportunizar a confecção e modelagem da massinha caseira tingida com anilina. - Possibilitar colagens com figuras \textbf{recortadas} de revistas, fotos, pedaços de tecidos de diferentes \textbf{texturas}, etc. - Organizar e possibilitar a construção de escultura com legumes, gravetos, folhas secas, blocos, copos \textbf{plásticos}, \textbf{embalagens} de papelão etc. - Possibilitar vivências de exploração com efeitos de luz e sombra sobre objetos ou espaços, com uso de velas, lanternas, data show e retroprojetor. - Organizar situações para que pintem usando diferentes suportes ( papéis, panos, telas, pedaços de metal ou acrílico, madeira, etc. ) e materiais ( aquarela, tinta guache, \textbf{tinta} feita com materiais da natureza, \textbf{lápis} de cor, canetas hidrográficas, esmalte de unhas e etc. ). - Possibilitar reconhecimento da diversidade de padrões de uso das cores em diferentes culturas e contextos de produção e que usem esse conhecimento para fazer suas criações no desenho, na \textbf{pintura}, etc. - Organizar e propiciar a utilização de objetos sonoros e instrumentos musicais em improvisações e composições. - Mediar a construção de instrumentos musicais de percussão, de sopro, de corda etc., com materiais alternativos. - Propiciar que \textbf{contem} histórias usando modulações de voz, objetos sonoros e instrumentos musicais. - Organizar situações para que \textbf{demonstrem} interesse por músicas de diferentes gêneros, estilos, épocas e culturas. - Possibilitar a participação na organização de cenário, a iluminação e o \textbf{som} para uma apresentação de teatro. - Oportunizar situações para a criação de formas planas e volumosas por meio da escultura, modelagem etc., e expressem opiniões sobre seu processo de produção. - Propiciar situações de construção de \textbf{brinquedos}, potes, cestos ou adornos inspirados no artesanato do campo, indígena ou de outras tradições culturais. - Possibilitar que construam casas ou castelos de cartas, de madeira, de panos e outros materiais. - Possibilitar a produção de dobraduras simples, bonecas de pano, de espiga de milho e etc. - Oportunizar situações para que possam desenhar com canetas hidrográficas em uma transparência e projetem na \textbf{parede} ou em uma tela ou lençol.

% matched lemmas: bater, bebê, brincadeira, brinquedo, contar, criança, demonstrar, embalagem, lápis, parede, pequeno, peso, pintura, plástico, pé, recortar, som, textura, tinta
\end{textsample}
